% ---------------------------------------------------------------
% IV. Deep Learning Architecture  (Member 1)
% ---------------------------------------------------------------

\section{Deep Learning Architecture and Core Pipeline}
\label{sec:dl_architecture}

\subsection{Data Split and Preprocessing}

We keep the cleaning of Phase~1: the 637 exact duplicates and the scans with no detected access point
are removed, leaving 19{,}227 training fingerprints. The official validation file (1{,}111 scans) is the
held-out \emph{test} set and is used only once, for the final results. To tune and early-stop the
networks we need a separate validation set, which we carve out of the training file: 15\% of the
\emph{reference positions} (2{,}795 scans) go to validation and the rest (16{,}432 scans) to training.
The split is grouped by position because each position was scanned about 20 times; a random split would
place near-identical fingerprints on both sides and make validation scores over-optimistic. The resulting
train and validation sets share no position, and all 13 building--floor classes appear in both.

Each fingerprint uses the RSSI of the access points observed at least once in the training part
($d = 459$). Undetected access points (coded $+100$) are set to $-110$~dBm, below the weakest real reading,
and every feature is mapped to $[0, 1]$ with a min--max scaler fitted on the training part only. Building
and floor labels are integer-encoded. The coordinates (UTM metres) are standardised with the training mean
and standard deviation, since raw values of the order of $4.86\times10^{6}$ would dominate the loss;
predictions are transformed back to metres before any error is computed.

\subsection{Multitask Network}

The network has a shared trunk of three fully connected blocks with 512, 256 and 128 units. Each block
applies a linear layer, layer normalisation and a GELU activation; the first two also apply dropout with
probability 0.2. The 128-dimensional output is a representation shared by three task heads: a linear
layer with 3 outputs (building logits), one with 5 outputs (floor logits) and one with 2 outputs
(standardised coordinates). The model has about $4.0\times10^{5}$ trainable parameters. Sharing the trunk
is the hard parameter sharing form of multitask learning described by Caruana~\cite{caruana1997multitask}:
the training signal of each task acts as an inductive bias for the others.

The training objective is the unweighted sum
\begin{equation}
\mathcal{L} = \mathcal{L}_{\text{building}} + \mathcal{L}_{\text{floor}} + \mathcal{L}_{\text{coords}},
\end{equation}
where the two classification losses are cross-entropies on the logits and $\mathcal{L}_{\text{coords}}$ is
the mean squared error on the standardised coordinates. Loss weighting is studied separately in
Section~VI. We train with Adam (learning rate $10^{-3}$, batch size 128) for at most 150 epochs.

\subsection{Model Selection}

Early stopping on the validation loss fails for this problem. After two epochs the validation floor
cross-entropy starts to rise, because the network becomes over-confident on a few hard scans, although
floor accuracy remains stable and the coordinate error keeps decreasing for many more epochs. Because the
classification losses are also about an order of magnitude larger than the coordinate loss, the summed
validation loss effectively monitors the floor alone. We therefore select the epoch that minimises the
evaluation score of the EvAAL/IPIN~2015 competition on this dataset~\cite{evaal2015},
\begin{equation}
s = \frac{1}{N}\sum_{i=1}^{N} \Big( \lVert \mathbf{p}_i - \hat{\mathbf{p}}_i \rVert_2
      + 50\,\mathbb{1}[\hat{b}_i \neq b_i] + 4\,\lvert \hat{f}_i - f_i \rvert \Big),
\end{equation}
computed on the validation set, where $\mathbf{p}_i$, $b_i$ and $f_i$ are the true position, building and
floor of scan $i$. Training stops after 20 epochs without improvement and the best checkpoint is restored.

\subsection{Single-Task Baselines}

To measure the effect of sharing, we train three single-task networks with exactly the same trunk,
preprocessing, optimiser, split and stopping rule, each with only one head: building only, floor only and
coordinates only. Their selection score uses only the term of $s$ that they predict. Every model is
trained with three random seeds (42, 7, 123), and we report the mean and standard deviation on the test set.

\subsection{Results}
Table~\ref{tab:dl_results} reports test-set results averaged over three seeds. Both network types clearly
improve positioning over the Phase~1 Random Forest (mean error 12.30~m, floor accuracy 89.92\%): the multitask
network reaches a mean error of 9.03~m and a floor accuracy of 91.21\%. Compared with the single-task networks,
however, sharing the trunk brings no measurable benefit under the unweighted loss: floor accuracy is equal
within the seed-to-seed spread, while the coordinate error (9.03 vs 8.59~m) and building accuracy (99.88 vs
99.97\%) are slightly worse. This is consistent with Caruana's observation that multitask learning is an
inductive bias that can either help or hurt~\cite{caruana1997multitask}. A plausible cause is the imbalance of
loss scales: on the validation set the floor cross-entropy is roughly twenty times larger than the
coordinate loss, so the shared representation is driven mainly by the floor task. Section~VI examines loss
weighting as a remedy.

\begin{table}[!t]
\centering
\caption{Test-set results (mean $\pm$ standard deviation over three seeds)}
\label{tab:dl_results}
\begin{tabular}{lcc}
\toprule
 & \textbf{Multitask} & \textbf{Single-task} \\
\midrule
Building accuracy (\%) & $99.88 \pm 0.14$ & $99.97 \pm 0.05$ \\
Floor accuracy (\%) & $91.21 \pm 0.23$ & $91.51 \pm 0.68$ \\
Mean error (m) & $9.03 \pm 0.10$ & $8.59 \pm 0.11$ \\
Median error (m) & $7.17 \pm 0.14$ & $6.30 \pm 0.13$ \\
P90 error (m) & $17.70 \pm 0.27$ & $16.91 \pm 0.25$ \\
\bottomrule
\end{tabular}
\end{table}
